% ============================================================================
%  PRISM-4D Preprint v4.0 -- Reframed for Peer Review
%  Based on: PRISM4D_Preprint_v3-3.pdf (Serfaty, 2026)
%  Target: Bioinformatics / JCTC / Nature Methods
%  Changes from v3.3: Tone reframing, equation numbering, expanded methods,
%                       expanded references, structural corrections
% ============================================================================

\documentclass[11pt,twocolumn]{article}

% ---- Packages ----
\usepackage[utf8]{inputenc}
\usepackage[T1]{fontenc}
\usepackage{amsmath,amssymb}
\usepackage{graphicx}
\usepackage{booktabs}
\usepackage{hyperref}
\usepackage[margin=1in]{geometry}
\usepackage{natbib}
\usepackage{xcolor}
\usepackage{algorithm}
\usepackage{algorithmic}
\usepackage{siunitx}
\usepackage{multirow}
\usepackage{caption}

% ---- Metadata ----
\title{PRISM-4D: Spike-Event-Driven Detection of Cryptic Binding Sites via Physics Simulation}

\author{Ididia Serfaty \\
  Delfictus IO LLC \\
  \texttt{is@delfictus.com}}

\date{February 2026}

\begin{document}
\maketitle

% ============================================================================
%  ABSTRACT
% ============================================================================
\begin{abstract}

\textbf{Background:}
Cryptic binding sites---transient pockets absent in static crystal structures---represent
high-value drug targets inaccessible to conventional computational screening. Existing
methods rely on surface geometry (fpocket, SiteMap) or machine learning trained on
crystallographic features (P2Rank, DeepSite), limiting their capacity to detect
conformationally gated pockets.

\textbf{Methods:}
We introduce PRISM-4D (Physics-driven Recognition of Intrinsic Sites through Molecular 4D
dynamics), a molecular dynamics platform that detects cryptic binding sites through explicit
physics simulation. PRISM-4D couples AMBER ff14SB force field molecular dynamics with a
5-phase cryo-thermal hysteresis protocol (50--300--50\,K, 55{,}000 steps/stream),
multi-wavelength UV excitation spectroscopy simulation (280/274/258/254/211\,nm), and a
leaky integrate-and-fire (LIF) neuromorphic network to identify cooperative dewetting events
indicative of druggable pocket formation. Detected spike events are clustered via a 4-stage
GPU pipeline: RT-DBSCAN via OptiX BVH traversal, watershed segmentation, Eikonal BFS, and
intensity$^2$-weighted centroid tracking (Spike-Native Density Clustering, SNDC).

\textbf{Results:}
On a 12-protein benchmark spanning orthosteric, allosteric, and cryptic binding sites,
PRISM-4D detects all 11 ground-truth targets within 10\,\AA{} Distance from Centroid to
Centroid (DCC), with 4/11 (36.4\%) at $<$5\,\AA{} and 8/11 (72.7\%) at $<$8\,\AA{}.
On the SIRP$\alpha$ immune checkpoint protein, PRISM-4D detects the experimentally
validated WYF cryptic pocket (7.1\,\AA{} DCC) as its top-ranked site from the closed apo
structure. On $\beta$-catenin, PRISM-4D identifies a computationally predicted cryptic
pocket (Site~6) featuring a dual-cysteine pair (CYS240/CYS278) with quality 0.965 and
centroid displacement $<$0.06\,\AA{} across 5 stochastic seeds.

\textbf{Conclusions:}
These results demonstrate that spike-event-driven physics simulation can access drug targets
fundamentally invisible to static-structure methods. The platform runs on a single
consumer-grade GPU with no external software dependencies.
\end{abstract}

\textbf{Keywords:} cryptic binding sites, molecular dynamics, drug discovery, spike detection,
pocket detection, cryo-thermal protocol, UV perturbation spectroscopy, AMBER force field,
RT-core acceleration, consumer GPU

% ============================================================================
%  1. INTRODUCTION
% ============================================================================
\section{Introduction}

The identification of druggable binding sites on protein surfaces is a fundamental step in
structure-based drug design. While orthosteric binding pockets are readily detectable from
static crystal structures, a growing body of evidence demonstrates that many therapeutically
relevant binding sites are cryptic: they exist only transiently, requiring conformational
rearrangement to become accessible to small molecules
\citep{cimermancic2016,vajda2018}. Cryptic pockets represent an enormous untapped
reservoir of drug targets, particularly for proteins historically classified as
``undruggable.''

Current computational approaches fall into two categories. Geometric methods (fpocket
\citep{leguilloux2009}, LIGSITE \citep{hendlich1997}, SiteMap \citep{halgren2009}) analyze the
static molecular surface using Voronoi tessellation or grid-based algorithms. Machine learning
methods (P2Rank \citep{krivak2018}, DeepSite \citep{jimenez2017}, PocketMiner
\citep{meller2023}) train on crystallographic features from databases such as sc-PDB. Both
share a fundamental limitation: they operate on a single static conformation and cannot detect
pockets requiring protein dynamics to open. FTMap \citep{kozakov2015} and mixed-solvent MD
\citep{ghanakota2016} partially address this through probe binding sampling, but remain
limited by short simulation timescales and lack explicit physics-driven detection criteria.

We present PRISM-4D (Physics-driven Recognition of Intrinsic Sites through Molecular 4D
dynamics), which replaces geometric inference and statistical learning with direct physics
simulation while running entirely on a single consumer-grade NVIDIA GPU with no external
software dependencies.

% ============================================================================
%  2. METHODS
% ============================================================================
\section{Methods}

\subsection{System Architecture}

PRISM-4D integrates four computational stages into a unified GPU-accelerated pipeline
implemented in Rust with CUDA kernels for performance-critical operations (Figure~1). The
compiled binary requires only the NVIDIA CUDA runtime and operates fully air-gapped with no
internet connectivity post-installation.

% ---- 2.1 Force Field and Integration ----
\subsection{Force Field and Integration}

Molecular dynamics simulations employ the AMBER ff14SB force field \citep{maier2015} with
the following protocol:
\begin{itemize}
  \item Integration: velocity Verlet with Langevin thermostat \citep{izaguirre2001}
  \item Timestep: 1\,fs (TODO: verify from code)
  \item Nonbonded cutoff: (TODO: extract from code)
  \item 1-4 scaling: LJ 0.5, electrostatic 5/6 (standard AMBER)
  \item Bond constraints: SHAKE \citep{ryckaert1977} on all X--H bonds
  \item Water model: implicit via hydrophobic exclusion inference (Section~\ref{sec:water})
\end{itemize}

% ---- 2.2 Cryo-Thermal Hysteresis Protocol ----
\subsection{Cryo-Thermal Hysteresis Protocol}
\label{sec:cryo}

The simulation executes a 5-phase thermal hysteresis cycle (Figure~2) selected by the
\texttt{--fast --hysteresis} flags:

\begin{enumerate}
  \item \textbf{Cold hold} (14{,}000 steps, 50\,K): Frozen baseline for clean UV spike detection
  \item \textbf{Ramp up} (6{,}000 steps, 50$\to$300\,K): Forward conformational transitions
  \item \textbf{Warm hold} (15{,}000 steps, 300\,K): Equilibrium sampling
  \item \textbf{Ramp down} (6{,}000 steps, 300$\to$50\,K): Reverse transitions; hysteresis measurement
  \item \textbf{Cold return} (14{,}000 steps, 50\,K): Final frozen state verification
\end{enumerate}

Total: 55{,}000 steps per stream. UV bursts fire every 250 steps, cycling through
wavelengths [280, 274, 258, 211]\,nm with 300-step dwell per wavelength (full cycle: 1{,}200
steps), at 42\,kcal/mol per burst.

% ---- 2.3 UV Pump-Probe Spectroscopy ----
\subsection{In Silico UV Pump-Probe Spectroscopy}

UV ``pulses'' deposit energy into aromatic residues according to experimentally measured
molar extinction coefficients. The energy deposition per chromophore is:
\begin{equation}
  E_{\text{dep}} = E_{\text{photon}} \times \sigma(\lambda) \times F \times \eta
  \label{eq:energy_dep}
\end{equation}
where the absorption cross-section $\sigma(\lambda)$ is converted from molar extinction
coefficient $\varepsilon$ via:
\begin{equation}
  \sigma = \varepsilon \times 3.823 \times 10^{-5} \; \text{\AA}^2
  \label{eq:sigma_conversion}
\end{equation}
$F$ is the fluence (0.024 photons/\AA$^2$), and $\eta$ is the quantum yield for thermal
dissipation. Chromophore parameters are summarized in Table~\ref{tab:chromophores}.

\begin{table}[h]
\centering
\caption{UV chromophore photophysical parameters.}
\label{tab:chromophores}
\begin{tabular}{lcccc}
\toprule
Chromophore & $\lambda_{\max}$ (nm) & $\varepsilon$ (M$^{-1}$cm$^{-1}$) & $\eta$ \\
\midrule
Trp & 280 & 5{,}600 & 0.83 \\
Tyr & 274 & 1{,}490 & 0.86 \\
Phe & 258 & 197 & 0.94 \\
S--S & 211 & 300 & 0.95 \\
Benzene (virtual) & 254 & 204 & 0.71 \\
\bottomrule
\end{tabular}
\end{table}

% ---- 2.4 Three-Channel Spike Detection ----
\subsection{Three-Channel Spike Detection Network}

Each voxel is monitored by three independent neuromorphic channels:

\textbf{Channel 1 (UV, source=1):} Active during UV bursts; detects aromatic excitation
within 4.0\,\AA{} with cooperative enhancement. Records wavelength, aromatic type, and
vibrational energy.

\textbf{Channel 2 (LIF, source=2):} Active when UV bursts are off; a leaky integrate-and-fire
neuron ($\tau_{\text{eff}}$\,=\,0.15\,ps, threshold\,=\,0.5, refractory\,=\,250 steps)
integrating bidirectional water density change and exclusion field signals. Detects
cooperative dewetting events.

\textbf{Channel 3 (EFP, source=3):} Electrostatic flux probe, always active when $\geq$2
charged atoms are nearby. Computes Coulomb potential with Warshel distance-dependent
dielectric:
\begin{equation}
  \varepsilon_r = \max(4r, 4)
  \label{eq:warshel}
\end{equation}
Emits typed CATION ($\varphi > 0$) or ANION ($\varphi < 0$) spikes.

Each spike event carries 96 bytes of metadata: position, intensity, source channel,
wavelength, aromatic type, water density, vibrational energy, and 8 nearby residue IDs.

% ---- 2.5 SNDC Clustering ----
\subsection{Spike-Native Density Clustering (SNDC v9)}

\subsubsection{RT-DBSCAN (Stage 3a)}
Spike positions are encoded as spheres in a bounding volume hierarchy (BVH) constructed by
NVIDIA OptiX RT cores. Epsilon-neighborhood queries become ray-sphere intersection tests,
reducing neighbor search from $O(N^2)$ to $O(N)$. The epsilon parameter is adaptively
computed:
\begin{equation}
  \epsilon = 3.0 \times \left(\frac{500}{N_{\text{atoms}}}\right)^{1/3}, \quad
  \epsilon \in [1.2, 3.0] \; \text{\AA}
  \label{eq:adaptive_eps}
\end{equation}

\subsubsection{Watershed Segmentation (Stage 3b)}
When a DBSCAN cluster spans multiple binding subsites, watershed segmentation identifies
density peaks within the cluster and partitions voxels by steepest-descent flow to each peak.

\subsubsection{Eikonal BFS (Stage 3c)}
Distance-field propagation from each watershed-separated cluster boundary via the Fast
Marching Method, enabling pocket geometry characterization: depth, entry channel width,
enclosure fraction, and solvent-accessible volume.

\subsubsection{Peak Centroid Tracking (Stage 3d)}
Intensity$^2$-weighted centroid computation:
\begin{equation}
  \mathbf{C} = \frac{\sum_i I_i^2 \cdot \mathbf{r}_i}{\sum_i I_i^2}
  \label{eq:centroid}
\end{equation}
The squared weighting concentrates the centroid on the thermodynamic hotspot rather than the
volumetric center, providing 3--15\,\AA{} improvement over geometric centroids for oversized
or irregular pockets.

% ---- 2.6 Benzene Cosolvent Probing ----
\subsection{UV-Activated Benzene Cosolvent Probing}

Virtual benzene probes are injected at surface-exposed hydrophobic aliphatic residues (LEU,
ILE, VAL, ALA, MET), using experimentally measured benzene photophysical properties
($\lambda_{\max}$\,=\,254\,nm, $\varepsilon$\,=\,204\,M$^{-1}$cm$^{-1}$, $\eta$\,=\,0.71;
\citealt{pantos1978,cundall1972}). The 254\,nm wavelength is integrated into the
frequency-hopping UV burst schedule alongside native aromatic channels (280/274/258/211\,nm).

% ---- 2.7 Water Inference ----
\subsection{Hydrophobic Exclusion Water Inference}
\label{sec:water}

Rather than employing explicit solvent simulation, PRISM-4D infers water density from the
hydrophobic exclusion field---the ``negative space'' defined by the protein's hydrophobic
surface. The exclusion field is computed via GPU-accelerated geometric kernels operating in
Fourier space.

% ---- 2.8 Benchmark Protocol ----
\subsection{Benchmark Protocol}

All benchmark runs used the command:
\begin{verbatim}
nhs_rt_full -t <topology>.json
  -o <output_dir> --fast --hysteresis
  --multi-stream 8 -v
\end{verbatim}
with \texttt{RUST\_LOG=info}. Topology preparation used \texttt{scripts/prism-prep} with
AMBER ff14SB force fields (requires OpenMM for \texttt{--use-amber}). Detection accuracy was
measured by Distance from Centroid to Centroid (DCC): the Euclidean distance between the
predicted site centroid and the crystallographic ligand centroid after Kabsch superposition
on backbone C$\alpha$ atoms. Each run uses 8 parallel independent streams (unique PRNG seeds:
$42 + i \times 12345$ for stream $i$), with consensus merge requiring a site to appear in
$\geq$4/8 streams (50\% threshold) within 10\,\AA{} spatial tolerance.

All benchmarks were executed on a single NVIDIA GeForce RTX 5080 (16\,GB GDDR7, Blackwell
architecture, compute capability 12.0).

% ============================================================================
%  3. RESULTS
% ============================================================================
\section{Results}

\subsection{Twelve-Protein Benchmark}

We evaluated PRISM-4D on 12 proteins spanning orthosteric, allosteric, and cryptic binding
sites (Table~\ref{tab:benchmark}). $\beta$-catenin (2GL7) is excluded from DCC accuracy
calculations as a computational prediction with no ground-truth ligand.

% TODO: Table 1 with error bars (mean +/- std from 8 streams)
% TODO: Add DCC cumulative distribution / ROC curve

All 11 ground-truth targets were detected within 10\,\AA{} DCC. Four targets achieved
excellent accuracy ($<$5\,\AA{}): BACE1 (3.6\,\AA{}), TEM1 $\beta$-lactamase (3.7\,\AA{}),
KRAS G12C (3.8\,\AA{}), and PTP1B (4.8\,\AA{}). Three marginal cases (9.5--9.8\,\AA{} DCC)
suggest refinement opportunities for multi-chain complexes and covalent cofactor geometries.

\subsection{Case Study 1: SIRP$\alpha$ WYF Cryptic Pocket}

Signal-regulatory protein alpha (SIRP$\alpha$) is a critical immune checkpoint in the
CD47--SIRP$\alpha$ axis. In December 2025, Diamond Light Source researchers discovered the WYF
cryptic pocket via an XChem fragment screening campaign (800+ fragments, 27 hits)
\citep{sirpa_xchem2025}. PRISM-4D detected this pocket as Site~1 (top-ranked, 7.1\,\AA{} DCC)
from the closed apo structure where the pocket is completely sealed. P2Rank 2.5 assigns 0.2\%
probability to the same pocket (Table~\ref{tab:sirpa}).

\subsection{Case Study 2: $\beta$-Catenin Predicted Cryptic Pocket}

$\beta$-catenin is widely considered ``undruggable'' due to its elongated armadillo repeat
surface. PRISM-4D identified a computationally predicted cryptic pocket---Site~6---featuring
TRP242 (aromatic anchor), CYS240/CYS278 (dual-cysteine pair), and LYS204 (electrostatic
anchor), scoring 0.965 quality and 0.844 druggability (Table~\ref{tab:bcatenin}). P2Rank
ranks this region 8/8 (0.5\% probability). \textit{This prediction requires experimental
validation.}

\subsection{Reproducibility Validation}

Five independent runs (different random seeds, same binary) demonstrate 0.06\,\AA{} maximum
pairwise centroid displacement for $\beta$-catenin Site~6 (Table~\ref{tab:repro}).
% TODO: Extend reproducibility analysis to all 12 targets

% ============================================================================
%  4. COMPARISON TO EXISTING METHODS
% ============================================================================
\section{Comparison to Existing Methods}

Table~\ref{tab:comparison} positions PRISM-4D against established pocket detection methods.
% TODO: Add head-to-head benchmark with PocketMiner and FTMap on same 12 proteins

% ============================================================================
%  5. DISCUSSION
% ============================================================================
\section{Discussion}

\subsection{Thermodynamic Event Detection}

PRISM-4D approaches pocket detection as a thermodynamic event detection problem. A druggable
pocket is a region where small molecules can be confined, water is excluded, and aromatic
residues create favorable interaction geometries. The LIF network detects exactly these
conditions through spike events at biophysically meaningful thresholds.

\subsection{Technical Contributions}

PRISM-4D introduces several techniques not previously applied to pocket detection
(Table~\ref{tab:innovations}), including in silico UV pump-probe spectroscopy, UV-activated
cosolvent probing, three-channel neuromorphic spike detection with temporal isolation, and
RT-core accelerated spatial clustering.

\subsection{The Dual-Cysteine Opportunity on $\beta$-Catenin}

The CYS240/CYS278 pair in $\beta$-catenin Site~6 represents a computationally predicted
covalent drug design opportunity requiring experimental validation. Dual-cysteine pockets are
rare ($<$200 human proteins with surface-accessible Cys pairs within 10\,\AA{}) and would
offer two-point covalent attachment for enhanced selectivity. The TRP242 aromatic anchor
creates a triangular Cys--Trp--Cys pharmacophore geometry.

\subsection{Hardware Accessibility}

The complete PRISM-4D platform runs on a single consumer-grade NVIDIA RTX 5080 GPU
(16\,GB GDDR7). The compiled binary requires only the NVIDIA CUDA runtime, operates without
network connectivity, and has no external software dependencies. This enables cryptic pocket
detection on hardware available to academic laboratories, startups, and researchers in
resource-limited settings.

\subsection{Limitations and Future Work}

\textbf{Benchmark size.} The current 12-protein benchmark is modest. Validation on the full
CryptoSite dataset \citep{cimermancic2016} ($\sim$55 proteins) and PocketMiner benchmark
\citep{meller2023} is ongoing.

\textbf{False positive rate.} The false discovery rate across all detected sites has not been
systematically quantified. PRISM-4D detects multiple sites per protein (e.g., 17 on
SIRP$\alpha$), and the proportion corresponding to true binding events requires systematic
evaluation.

\textbf{Negative controls.} The current benchmark lacks proteins with no known cryptic
pockets. Validation on established negative controls is needed to assess specificity.

\textbf{Parameter sensitivity.} Key parameters (LIF threshold, UV burst energy, epsilon
range) were empirically tuned. Systematic sensitivity analysis across diverse protein
families is planned.

\textbf{Multi-chain complexes.} Three marginal results (9.5--9.8\,\AA{} DCC) involve
multi-chain complexes or covalent cofactor geometries, suggesting algorithmic refinement
opportunities.

\textbf{GPU vendor dependency.} The current implementation requires NVIDIA GPUs with CUDA
support and RT cores (Turing architecture or newer). AMD and Intel GPU support is not
available.

\textbf{Temperature protocol.} The 50\,K cold phase employs temperatures below the glass
transition of solvated proteins ($\sim$180\,K), which may introduce non-physiological
conformational states. The extent to which these affect pocket detection accuracy warrants
investigation.

\textbf{Runtime.} Execution time (3--18 minutes per target) is substantially slower than
geometric methods (seconds) though faster than conventional MD approaches (hours to days).

% ============================================================================
%  6. CONCLUSION
% ============================================================================
\section{Conclusion}

PRISM-4D introduces a spike-event-driven approach to binding site detection based on physics
simulation. Its SNDC pipeline---RT-DBSCAN, watershed segmentation, Eikonal BFS, and peak
centroid tracking---detects all 11 ground-truth targets within 10\,\AA{} DCC on a 12-protein
benchmark, identifies the experimentally validated SIRP$\alpha$ WYF cryptic pocket as its
top-ranked site from the closed apo structure, and predicts a novel cryptic pocket on
$\beta$-catenin with sub-angstrom centroid reproducibility across stochastic seeds. The
platform achieves these results on consumer-grade hardware with no external dependencies.

% ============================================================================
%  7. DATA AND CODE AVAILABILITY
% ============================================================================
\section{Data and Code Availability}

PRISM-4D is developed by Delfictus IO LLC. The engine is implemented in Rust with CUDA GPU
acceleration (Blackwell architecture, compute capability $\geq$12.0). Benchmark topologies,
binding site JSON outputs, and PyMOL session files are deposited at [TODO: Zenodo DOI].
% TODO: Deposit data and provide DOI
% TODO: Code availability statement (open-source timeline or embargo)

% ============================================================================
%  REFERENCES
% ============================================================================
\bibliographystyle{plainnat}
% TODO: Create .bib file with 40+ references
% \bibliography{prism4d}

\begin{thebibliography}{99}

\bibitem[Bowman and Geissler(2012)]{bowman2012}
Bowman, G.R. and Geissler, P.L. (2012).
Equilibrium fluctuations of a single folded protein reveal a multitude of potential cryptic allosteric sites.
\textit{Proc. Natl. Acad. Sci.}, 109(29), 11681--11686.

\bibitem[Cimermancic et al.(2016)]{cimermancic2016}
Cimermancic, P., Weinkam, P., Rettenmaier, T.J., et al. (2016).
CryptoSite: Expanding the druggable proteome by characterization and prediction of cryptic binding sites.
\textit{J. Mol. Biol.}, 428(4), 709--719.

\bibitem[Cundall and Robinson(1972)]{cundall1972}
Cundall, R.B. and Robinson, D.A. (1972).
Primary photophysical processes in aromatic molecules.

\bibitem[Durrant and McCammon(2011)]{durrant2011}
Durrant, J.D. and McCammon, J.A. (2011).
Molecular dynamics simulations and drug discovery.
\textit{BMC Biol.}, 9(1), 71.

\bibitem[Ester et al.(1996)]{ester1996}
Ester, M., Kriegel, H.P., Sander, J., and Xu, X. (1996).
A density-based algorithm for discovering clusters in large spatial databases with noise.
\textit{KDD}, 226--231.

\bibitem[Ghanakota and Carlson(2016)]{ghanakota2016}
Ghanakota, P. and Carlson, H.A. (2016).
Driving structure-based drug discovery through cosolvent molecular dynamics.
\textit{J. Med. Chem.}, 59(23), 10383--10399.

\bibitem[Halgren(2009)]{halgren2009}
Halgren, T.A. (2009).
Identifying and characterizing binding sites and assessing druggability.
\textit{J. Chem. Inf. Model.}, 49(2), 377--389.

\bibitem[Hendlich et al.(1997)]{hendlich1997}
Hendlich, M., Rippmann, F., and Barnickel, G. (1997).
LIGSITE: automatic and efficient detection of potential small molecule-binding sites in proteins.
\textit{J. Mol. Graph. Model.}, 15(6), 359--363.

\bibitem[Izaguirre et al.(2001)]{izaguirre2001}
Izaguirre, J.A., Catarello, D.P., Wozniak, J.M., and Skeel, R.D. (2001).
Langevin stabilization of molecular dynamics.
\textit{J. Chem. Phys.}, 114(5), 2090--2098.

\bibitem[Jimenez et al.(2017)]{jimenez2017}
Jimenez, J., Doerr, S., Martinez-Rosell, G., Rose, A.S., and De Fabritiis, G. (2017).
DeepSite: protein-binding site predictor using 3D-convolutional neural networks.
\textit{Bioinformatics}, 33(19), 3036--3042.

\bibitem[Kozakov et al.(2015)]{kozakov2015}
Kozakov, D., Grove, L.E., Hall, D.R., et al. (2015).
The FTMap family of web servers for determining and characterizing ligand-binding hot spots.
\textit{Nat. Protoc.}, 10(5), 733--755.

\bibitem[Krivak and Hoksza(2018)]{krivak2018}
Krivak, R. and Hoksza, D. (2018).
P2Rank: machine learning based tool for rapid and accurate prediction of ligand binding sites.
\textit{J. Cheminform.}, 10(1), 39.

\bibitem[Lakowicz(2006)]{lakowicz2006}
Lakowicz, J.R. (2006).
\textit{Principles of Fluorescence Spectroscopy}, 3rd ed. Springer.

\bibitem[Le~Guilloux et al.(2009)]{leguilloux2009}
Le~Guilloux, V., Schmidtke, P., and Tuffery, P. (2009).
Fpocket: an open source platform for ligand pocket detection.
\textit{BMC Bioinform.}, 10(1), 168.

\bibitem[Maier et al.(2015)]{maier2015}
Maier, J.A., Martinez, C., Kasavajhala, K., et al. (2015).
ff14SB: Improving the accuracy of protein side chain and backbone parameters from ff99SB.
\textit{J. Chem. Theory Comput.}, 11(8), 3696--3713.

\bibitem[Meller et al.(2023)]{meller2023}
Meller, A., Ward, M., Borowsky, J., et al. (2023).
Predicting locations of cryptic pockets from single protein structures using the PocketMiner graph neural network.
\textit{Nat. Commun.}, 14(1), 1177.

\bibitem[Pace et al.(1995)]{pace1995}
Pace, C.N., Vajdos, F., Fee, L., Grimsley, G., and Gray, T. (1995).
How to measure and predict the molar absorption coefficient of a protein.
\textit{Protein Sci.}, 4(11), 2411--2423.

\bibitem[Pantos et al.(1978)]{pantos1978}
Pantos, E., Philis, J., and Bolovinos, A. (1978).
The extinction coefficient of benzene vapor in the region 4.6 to 36 eV.
\textit{J. Mol. Spectrosc.}, 72(1), 36--43.

\bibitem[Ryckaert et al.(1977)]{ryckaert1977}
Ryckaert, J.P., Ciccotti, G., and Berendsen, H.J.C. (1977).
Numerical integration of the Cartesian equations of motion of a system with constraints.
\textit{J. Comput. Phys.}, 23(3), 327--341.

\bibitem[Vajda et al.(2018)]{vajda2018}
Vajda, S., Beglov, D., Wakefield, A.E., et al. (2018).
Cryptic binding sites on proteins: definition, detection, and druggability.
\textit{Curr. Opin. Chem. Biol.}, 44, 1--8.

\bibitem[Vincent and Soille(1991)]{vincent1991}
Vincent, L. and Soille, P. (1991).
Watersheds in digital spaces: an efficient algorithm based on immersion simulations.
\textit{IEEE Trans. Pattern Anal. Mach. Intell.}, 13(6), 583--598.

\bibitem[Wetlaufer(1962)]{wetlaufer1962}
Wetlaufer, D.B. (1962).
Ultraviolet spectra of proteins and amino acids.
\textit{Advances in Protein Chemistry}, 17, 303--390.

\bibitem[SIRPa~XChem(2025)]{sirpa_xchem2025}
[Authors TBD] (2025).
Engineering SIRP$\alpha$ conformational plasticity to reveal a cryptic pocket suitable for structure-based drug design.
bioRxiv preprint, December 23, 2025. Diamond Light Source XChem fragment screening.

% TODO: Add remaining 15-20 references (see remediation tracker)

\end{thebibliography}

\end{document}
